\RequirePackage[2020-02-02]{latexrelease} % workaround for broken
\documentclass[10pt,a4paper]{article}
% ========================== DO NOT MODIFY =================================== %
\usepackage[a4paper,left=2.0cm,top=2.0cm,right=2.0cm,bottom=2.0cm]{geometry}
\usepackage{amssymb,amsmath} % If amsmath is required
\usepackage[implicit=false]{hyperref}
\usepackage[utf8]{inputenc}
\usepackage[T1]{fontenc}
\usepackage{times}
\usepackage{graphicx}
\usepackage{titling}
\usepackage{authblk}
\usepackage{multicol}
\usepackage{float}
\usepackage{enumitem}
\usepackage{booktabs}
\usepackage{xcolor}
\usepackage{hyperref}
\hypersetup{
    colorlinks=true,
    linkcolor=blue,
    filecolor=blue,      
    urlcolor=blue,
}
    
\usepackage{fancyhdr}
\fancyhf{}
\setlength{\headheight}{3em}
\pagestyle{fancy}
\renewcommand{\headrulewidth}{0pt} % remove lines as well
\fancyhead[L]{\footnotesize{24\textsuperscript{th} International Conference on Biomedical Applications of Electrical Impedance Tomography (EIT 2024), Hangzhou, China}}
\fancypagestyle{firstpage}{%
  \lhead{\footnotesize{24\textsuperscript{th} International Conference on Biomedical Applications of Electrical Impedance Tomography (EIT 2024),  Hangzhou, China}}
}

\usepackage[small,compact]{titlesec}
% Fix titlesec for the 2.10.1 version in ubuntu 16.04
%  https://bugs.launchpad.net/ubuntu/+source/texlive-extra/+bug/1574052
\makeatletter
\@ifpackagelater{titlesec}{2016/03/21}{%
 % Package titlesec is on version 2.10.2 or higher, nothing to do %
}{%
 % Check if package titlesec is on version 2.10.1 %
 \@ifpackagelater{titlesec}{2016/03/15}{%
  % Package titlesec on version 2.10.1, patch accordingly %
  \usepackage{etoolbox}%
  \makeatletter
  \patchcmd{\ttlh@hang}{\parindent\z@}{\parindent\z@\leavevmode}{}{}%
  \patchcmd{\ttlh@hang}{\noindent}{}{}{}%
  \makeatother
 }{%
  % Package titlsecon is on version 2.10.0 or lower, nothing to do %
 }%
}
\makeatother
\titlespacing*{\section}{0pt}{5pt}{5pt}  % If you need different spacing

\usepackage[numbers,sort&compress]{natbib}
\usepackage[small,bf,tableposition=top,figureposition=bottom,skip=2pt]{caption}


% limit space between floats and text
\setlength{\textfloatsep}{10pt plus 1.0pt minus 2.0pt}
% shrink the list environments
\setlist{noitemsep}
\setlist{nolistsep}

% Shrink the bibliography
  \let\oldthebibliography=\thebibliography
  \let\endoldthebibliography=\endthebibliography
  \renewenvironment{thebibliography}[1]{%
    \begin{oldthebibliography}{#1}%
      \setlength{\parskip}{0ex}%
      \setlength{\itemsep}{0ex}%
  }%
  {%
    \end{oldthebibliography}%
  }

% Format of the front matter
\renewcommand{\Affilfont}{\small}
\setlength{\affilsep}{1ex}
\renewcommand\Authands{, }
\setlength{\droptitle}{-2cm}
% ============================================================================ %




%------------------------------- FRONT MATTER -------------------------------- %
\title{Anatomically accurate torso mesh for EIT
\vspace{-0.5ex}} %remove vertical space
\author[1]{Bartłomiej Grychtol}
\author[1]{Andy Adler}
\affil[1]{Systems and Computer Engineering, Carleton University, Ottawa, Canada, \protect\href{mailto:adler@sce.carleton.ca}{\underline{adler@asce.carleton.ca}}}
\date{}
%----------------------------------------------------------------------------- %
%\protect\url


\begin{document}
\maketitle
\thispagestyle{firstpage}
\vspace{-1.5cm}
%\thispagestyle{empty}

\begin{multicols}{2}

\noindent {\bf Abstract:}
We describe an anatomically accurate torso mesh and set of tools to allow
arbitrary electrode placement. The work is made available through EIDORS and
on Zenodo.

\section{Introduction}
Anatomically-accurate finite element meshes are needed for EIT forward
and inverse simulations. While many such meshes are publically avaialble, 
EIT also requires the ability to specify and refine body surface electrodes.
To investigate novel EIT configurations (e.g.{} for 3D EIT) tools are needed
to support arbitrarily electrode positions. 

\section{Methods and Results}
A base mesh was taken from ???
Next, 

The mesh contains segmented airways and major vessels. The trachea and
major bronchi are segmented to the level of ???. The heart is segmented
into 
Vessels

Example image of the FEM mesh are shown in Fig.1. All software to generate the
component sub-images are make available at \cite{grychtol2024zenodo}.

\section{Discussion}
We hope this contribution helps the EIT and bioimpedance community with
to more easily model detailed anatomical features of the thoracic anatomy.
To facilitate access to this model, it has been made available as part of
EIDORS\cite{adler2006}, and also uploaded to zenodo\cite{grychtol2024zenodo}.






\footnotesize
% Instead of using thebibliography, you can also type the references
% using bibtex, and uncomment the following lines
%\bibliographystyle{compact}
%\bibliography{eit2018}
\begin{thebibliography}{}
\bibitem{adler2006}
Adler A, Lionheart WRB,
{\em Physiol Meas} 27:S25--S42, 2006.

\bibitem{grychtol2024zenodo}
Grychtol B, Adler A,
????? Zenodo link

\end{thebibliography}
\end{multicols}


\begin{figure}[H]
\centering
\includegraphics[width=.99\columnwidth,page=1]{figures.pdf}
\caption{%
Illustrations of the generated mesh with multiple electrode planes.
A) Full model with coloured organ regions, B) zoom into the full model
in the heard region, C) mesh in a slice through the ventricles.
}
\label{fig:loc}
\end{figure}


\end{document}
